\chapter{REVIEW OF RELATED LITERATURE}

\section{Brain Tumor Segmentation}

Approaches for segmenting brain tumors from medical imaging data includes traditional techniques like thresholding, region-growing, active contours and more recent deep learning methods like convolutional neural networks.

\subsection{Traditional Methods}

Traditional methods for brain tumor segmentation have relied on classical image processing techniques, such as thresholding, region-growing, and active contour models. Thresholding methods segment images by setting intensity values above or below a certain threshold to a specific value, effectively separating the tumor region from the surrounding tissues \parencite{gordillo2013state}. Region-growing algorithms, on the other hand, start from a seed point and iteratively expand the region by adding neighboring pixels that satisfy certain similarity criteria \parencite{melouah2018overview}.  Active contour models, also known as snakes, deform an initial contour towards the boundaries of the tumor region based on energy minimization principles \parencite{caselles1997geodesic}.

While these traditional methods have been widely used, they often face limitations in handling complex tumor shapes and intensities. For instance, thresholding techniques struggle to differentiate tumors from surrounding tissues with similar intensities, and they are susceptible to noise and artifacts \parencite{cai2018quantitative}. Region-growing methods can also be sensitive to initialization and may fail to accurately segment tumors with irregular shapes or heterogeneous intensities \parencite{gordillo2013state}. Active contour models, while capable of handling complex shapes, require careful parameter tuning and can be computationally intensive, limiting their practical applicability \parencite{pierre2021segmentation}.

\subsection{Deep Learning Methods}

The advent of deep learning, particularly Deep Convolutional Neural Networks (DCNNs), has revolutionized the field of medical image analysis, including brain tumor segmentation. DCNNs have demonstrated remarkable ability in capturing complex spatial relationships and learning hierarchical features from medical images, enabling more accurate and robust segmentation of brain tumors \parencite{litjens2017survey}.

One of the most widely used DCNN architectures for brain tumor segmentation is the U-Net \parencite{ronneberger2015u} a fully convolutional network that employs skip connections to preserve spatial information and capture both local and global features. Other popular architectures include SegNet \parencite{badrinarayanan2017segnet}, which uses an encoder-decoder structure with pooling indices for efficient upsampling, and DeepLabv3+ \parencite{chen2018encoder}, which incorporates atrous convolutions and pyramid pooling for improved segmentation performance.

These deep learning-based methods have consistently outperformed traditional techniques in brain tumor segmentation tasks, as they can effectively handle the complex and diverse appearances of brain tumors, as well as the presence of noise and artifacts \parencite{gupta2021deep}. However, the success of these methods heavily relies on the availability of large, high-quality datasets for training, which can be challenging to obtain in the medical domain.

\subsection{Challenges in Brain Tumor Segmentation}

Despite the advancements in deep learning-based brain tumor segmentation, several challenges remain:

\begin{enumerate}
    \item \textbf{Data scarcity and class imbalance:} Acquiring large, well-annotated datasets for brain tumor segmentation can be a significant challenge due to the time and expertise required for manual annotation \parencite{ullah2023enhancing}. Additionally, the presence of a large number of non-tumor voxels can lead to class imbalance, hindering the training of deep learning models \parencite{isensee2018brain}.
    \item \textbf{Tumor heterogeneity:} Brain tumors can exhibit diverse shapes, sizes, and intensities, making it difficult to generalize segmentation models to different types of tumors \parencite{ranjbarzadeh2021brain}. This heterogeneity can arise from various factors, such as different tumor grades, molecular subtypes, and the presence of associated pathologies like edema or necrosis \parencite{bakas2018identifying}.
    \item \textbf{Computational complexity:} Deep learning models, particularly those with complex architectures and high-resolution input images, can be computationally expensive to train and deploy \parencite{havaei2017brain}. This can pose challenges in clinical settings where computational resources may be limited, and real-time performance is crucial for timely diagnosis and treatment planning.
\end{enumerate}

\section{Attention Mechanisms in Deep Learning}

A powerful concept that enables deep learning models to selectively focus on the most relevant parts of input data and capture long-range dependencies, with applications in medical image segmentation.

\subsection{Types of Attention Mechanisms}

Attention mechanisms have emerged as a powerful tool in deep learning, allowing models to selectively focus on the most relevant parts of the input data and capture long-range dependencies. Several types of attention mechanisms have been developed, including:

\begin{enumerate}
    \item \textbf{Spatial attention:} Spatial attention mechanisms enable models to attend to specific regions of an image or feature map, allowing them to focus on the most relevant spatial locations \parencite{jetley2018learn}. This can be particularly useful in medical image analysis, where the regions of interest (e.g., tumors) may be localized within the image.
    \item \textbf{Channel attention:} Channel attention mechanisms allow models to learn the importance of different channels or features within a feature map \parencite{hu2018squeeze}. By selectively emphasizing or suppressing specific channels, the model can better capture the most discriminative features for the task at hand.
    \item \textbf{Self-attention:} Self-attention mechanisms, such as the Transformer architecture \parencite{vaswani2017attention}, enable models to attend to relationships between different elements within a sequence or feature map. This allows the model to capture long-range dependencies and learn more expressive representations, which can be beneficial for tasks like image segmentation, where global context is important.
\end{enumerate}

\subsection{Applications in Medical Image Segmentation}

Attention mechanisms have been successfully applied to various medical image segmentation tasks, including brain tumor segmentation. For instance, the study by \textcite{sinha2020multi} proposed an attention-based U-Net architecture that incorporates spatial and channel attention modules, demonstrating improved segmentation performance for brain tumors compared to the vanilla U-Net.

This study has shown that attention mechanisms can effectively guide deep learning models to focus on the most relevant features and spatial locations, leading to improved segmentation accuracy and robustness.

\section{Attention Synergy Network (AS-Net)}

A novel deep learning architecture designed for image segmentation tasks that combines spatial and channel attention mechanisms in a synergistic manner, demonstrated for skin lesion segmentation.

\subsection{Architecture and Methodology}

The Attention Synergy Network (AS-Net) \parencite{hu2022net} is a novel deep learning architecture designed for image segmentation tasks, particularly skin lesion segmentation. The AS-Net architecture combines spatial and channel attention mechanisms in a synergistic manner to enhance the model's ability to capture both local and global features.

The AS-Net model consists of an encoder network, which extracts feature maps from the input image, and a decoder network, which generates the segmentation output. The key components of the AS-Net architecture are:

\begin{enumerate}
    \item \textbf{Spatial attention module:} This module applies spatial attention to the feature maps produced by the encoder, allowing the model to focus on the most relevant spatial locations and suppress irrelevant regions.
    \item \textbf{Channel attention module:} This module applies channel attention to the feature maps, enabling the model to selectively emphasize or suppress specific channels or features based on their importance for the segmentation task.
    \item \textbf{Synergy module:} This module combines the outputs of the spatial and channel attention modules using a fusion mechanism, effectively integrating the spatial and channel-wise attention information to produce more informative feature representations.
\end{enumerate}

The authors of AS-Net demonstrated its effectiveness in skin lesion segmentation, achieving state-of-the-art performance on several benchmark datasets \parencite{hu2022net}. However, they acknowledged limitations in terms of computational efficiency and generalizability due to the use of the computationally heavy VGG16 encoder architecture.

\subsection{Performance in Skin Lesion Segmentation}

While the AS-Net model achieved impressive results in skin lesion segmentation, the authors identified several limitations that could hinder its practical applicability in other domains or real-world scenarios:

\begin{enumerate}
    \item \textbf{Computational efficiency:} The VGG16 encoder used in the original AS-Net architecture is computationally intensive, which can lead to slow inference times and higher resource requirements. This can be a significant drawback in medical imaging applications, where real-time performance and efficient deployment are often crucial.
    \item \textbf{Generalizability:} The AS-Net model was primarily developed and evaluated for skin lesion segmentation, and its generalizability to other types of medical images, such as MRI scans for brain tumor segmentation, has not been extensively studied. Different imaging modalities and anatomical structures may require adaptations or modifications to the model architecture and training strategies.
\end{enumerate}

To address these limitations, the authors suggested exploring alternative encoder architectures, particularly lightweight models, that could potentially improve computational efficiency without compromising segmentation performance. Additionally, they emphasized the need to evaluate the model's generalizability across different medical image segmentation tasks and datasets.

\section{Lightweight Encoder Architectures}

Computationally efficient encoder architectures like MobileNetV3 and EfficientNetV2 that balance accuracy and efficiency, explored for medical image segmentation tasks.

\subsection{Importance of Lightweight Models}

In the field of medical image analysis, there is a growing need for computationally efficient deep learning models that can perform segmentation tasks with high accuracy while minimizing computational resources and processing time. This need arises from several practical considerations in medical imaging applications:

\begin{enumerate}
    \item \textbf{Real-time performance:} In certain clinical scenarios, such as intraoperative guidance or real-time monitoring, segmentation models need to process and analyze images rapidly to provide timely feedback and support decision-making processes \parencite{ullah2023enhancing}. Computationally intensive models can introduce significant latency, hindering their practical utility in such time-sensitive applications.
    \item \textbf{Deployment on edge devices:} With the increasing adoption of telemedicine and remote patient monitoring, there is a growing demand for deploying medical image analysis models on resource-constrained devices, such as smartphones, tablets, or embedded systems \parencite{hirschorn2014use}. These edge devices often have limited computational power and memory, necessitating the use of lightweight and efficient models that can perform inference on-device without compromising performance.
    \item \textbf{Scalability and cost-effectiveness:} In large-scale healthcare systems or research facilities, deploying computationally intensive models can result in significant infrastructure costs and energy consumption \parencite{canziani2016analysis}. Lightweight models can help reduce these costs and improve the scalability of medical image analysis solutions, enabling wider adoption and accessibility.
\end{enumerate}

By exploring lightweight encoder architectures, the proposed study aims to address these practical considerations and enhance the computational efficiency of the AS-Net model for MRI-based brain tumor segmentation, potentially enabling real-time performance, on-device deployment, and cost-effective scaling in clinical settings. 


\subsection{Examples of Lightweight Encoders}

Two prominent examples of lightweight encoder architectures that have gained significant attention in recent years are MobileNetV3 and EfficientNetV2. These architectures have been specifically designed to strike a balance between computational efficiency and accuracy, making them suitable candidates for the proposed study.

\textbf{MobileNetV3} is a highly efficient convolutional neural network architecture developed by Google. It builds upon the success of previous MobileNet versions and introduces several improvements, such as improved efficiency through efficient network design and hardware-aware optimization techniques \parencite{howard2019searching}. MobileNetV3 employs a combination of depth-wise separable convolutions and inverted residual blocks, which significantly reduces the computational cost and model size compared to traditional convolutional architectures.

\textbf{EfficientNetV2}, on the other hand, is a family of convolutional neural network architectures designed by Google Brain, with a focus on achieving superior performance while maintaining computational efficiency. EfficientNetV2 models are based on a novel compound scaling method that uniformly scales the depth, width, and resolution of the network, enabling efficient scaling and improved accuracy-efficiency trade-offs \parencite{tan2021efficientnetv2}.

In the domain of medical image analysis, EfficientNetV2 has shown promising results in various segmentation tasks. For example, the study by \textcite{lin2024brain} integrated EfficientNetV2 as the encoder in a U-Net-based architecture for lung segmentation, demonstrating improved performance compared to traditional encoders like ResNet while maintaining computational efficiency.

Both MobileNetV3 and EfficientNetV2 have been extensively studied and evaluated across various computer vision tasks, demonstrating their ability to achieve high accuracy while significantly reducing computational complexity and memory footprint. However, their specific performance in the context of MRI-based brain tumor segmentation, particularly when integrated into the AS-Net architecture, has not been thoroughly investigated.

By comparing the performance of these lightweight encoder architectures within the AS-Net framework, the proposed study aims to provide valuable insights into the trade-offs between segmentation accuracy, computational efficiency, and generalizability. Ultimately, the findings of this research can contribute to the development of more practical and clinically-relevant brain tumor segmentation solutions that can be effectively deployed in real-world healthcare settings.
